\documentclass[11pt]{article}
\usepackage[a4paper,margin=25mm]{geometry}
\usepackage{amsmath,amssymb,amsthm,mathtools}
\usepackage[hidelinks]{hyperref}
\newtheorem{theorem}{Theorem}
\newtheorem{lemma}{Lemma}
\newcommand{\dd}{\,d}
\title{Original theta pressure and the sign of the arithmetic correction\\
Independent receiving derivation with the physical equilibrium constant}
\author{Split-Zero research continuation}
\date{19 September 2026}
\begin{document}
\maketitle

This is a successor proof to the accepted signed arithmetic calculation,
not an alteration of its sealed reader. The incoming sign argument is
rederived below. In particular the equilibrium logarithmic moment is
derived directly from the original EIQ measure, so its closed form is
not an additional assumption. All finite inequalities are rational;
\texttt{check\_theta\_sign\_exact.py} implements precisely their finite sums.
The classical theta and special-function inputs are cited at use.

\section{Original source, original coefficient, and result}
Fix an actual off-line quartet of zeros of $g=2\xi$, with complete common
order $m_0\geq1$ and representatives
\begin{equation}\tag{TS1}
 \rho=\tfrac12+\delta+i\gamma,\qquad 0<\delta<\tfrac12,\quad\gamma>2,
 \qquad
 h(s)=\prod_{z\in\{\rho,\bar\rho,1-\rho,1-\bar\rho\}}(s-z)^{m_0}.
\end{equation}
The word actual means that $g(\rho)=0$ and that these are the stated
zeros of the original numerator. No existence of such a quartet is
asserted here. Define, with its original mass,
\begin{equation}\tag{TS2}
 w_h(y)=\frac{|g(1/2+iy)/h(1/2+iy)|^2}{2\pi},\quad
 M_h(\theta)=\int_{\mathbb R}e^{\theta y}w_h(y)\dd y,\quad
 \alpha_*=\frac\pi2,\quad
 I_h=\int_0^{\alpha_*}\left(\frac{M_h'(\theta)}{M_h(\theta)}\right)^2\dd\theta.
\end{equation}
The already proved original-source envelope is
$e^{\alpha_*|y|}w_h(y)\leq C_h(1+|y|)^{-p_h}$ with
$p_h=8m_0-9/2>3$; its complete provider is recorded in
\cite{SBV}. Thus dominated convergence gives $M_h\in C^2$
on the closed tilt interval, with positive values there. The mass is
$M_h(0)$ and is not replaced by one.

Let $\rho_{a,b}$ denote the original equilibrium measure for
$V_{a,b}(x)=b\sqrt x-a\log x$ on $(0,\infty)$, and put
\begin{equation}\tag{TS3}
 \mathcal L(a,b)=\int\log x\dd\rho_{a,b}(x),\qquad
 \mathcal C_\Gamma=2\mathcal L(2,\pi)-8\log2+4.
\end{equation}
Section~\ref{eqphysical} retains its physical elliptic integrals,
constructs this measure, and computes the moment. The symbol
$\mathcal C_\Gamma$ refers to this return coefficient, not a constant
in a pointwise upper bound for the Gamma density.

\begin{theorem}\label{th:sign}
For every actual quartet in (TS1), at every complete multiplicity,
\begin{equation}\tag{TS4}
 \frac{M_h(\pi/2)}{M_h(0)}>61>e^{41/10},\qquad
 I_h>\frac{1681}{50\pi},\qquad 0<\mathcal C_\Gamma<\frac23.
\end{equation}
For $m_0=1$ the already evaluated coefficient therefore satisfies
\begin{equation}\tag{TS5}
 d_0(h)=-2\mathcal C_\Gamma+\frac{\log2}{\pi}I_h
 >\kappa_*:=\frac{7370383}{7260000}>1.
\end{equation}
The pressure bound (TS4) applies also when $m_0\geq2$.
At those fixed higher multiplicities the asymptotic coefficient is
$-(4m_0-2)\mathcal C_\Gamma<0$, because the original $k^2/q$ tends to zero.
\end{theorem}

\section{Positive theta kernel with the factor two retained}
For $x>0$ put
\[
 \vartheta(x)=\sum_{n\in\mathbb Z}e^{-\pi n^2x^2},\quad
 D=-x\partial_x,\quad (Jf)(x)=x^{-1}f(1/x).
\]
The theta transformation $\vartheta=J\vartheta$ is the specialization
of the classical modular identity, DLMF 20.7.32 at $z=0$,
$\tau=ix^2$ \cite{theta}. Differentiating a term with $a=\pi n^2$
gives $D e^{-ax^2}=2ax^2e^{-ax^2}$ and
$D^2e^{-ax^2}=(-4ax^2+4a^2x^4)e^{-ax^2}$. Therefore
\begin{equation}\tag{TS6}
 T(x)=D(D-1)\vartheta(x)
 =2\sum_{n\geq1}(4\pi^2n^4x^4-6\pi n^2x^2)e^{-\pi n^2x^2}.
\end{equation}
The series and all its derivatives converge uniformly on compact
subsets of $x>0$. Direct differentiation of $Jf$ gives
$DJ=J(1-D)$, hence $D(D-1)J=JD(D-1)$ and $T=JT$.
Consequently $K(t)=e^{t/2}T(e^t)$ is even. For $t\geq0$ every term
of (TS6) is strictly positive, since $2\pi n^2e^{2t}-3>0$.
It follows that $K(t)>0$ for every real $t$. For $x\geq1$, the
Gaussian series in (TS6) is bounded by $C(1+x^4)e^{-\pi x^2}$;
the same type of bound holds after any fixed number of logarithmic
derivatives. Reflection gives decay faster than every exponential
at the other end. This proves all Mellin and Fourier integrability
statements used next.

For $\Re s>1$ each Mellin integral is absolutely integrable and its
sum is absolutely convergent. Substitution $u=\pi n^2x^2$ gives
\[
 \int_0^\infty(4\pi^2n^4x^4-6\pi n^2x^2)e^{-\pi n^2x^2}x^s\frac{dx}{x}
 =\frac12\pi^{-s/2}n^{-s}s(s-1)\Gamma(s/2).
\]
Here $4\Gamma(s/2+2)-6\Gamma(s/2+1)=s(s-1)\Gamma(s/2)$.
The literal factor two in (TS6) yields
\begin{equation}\tag{TS7}
 \int_0^\infty T(x)x^s\frac{dx}{x}
 =s(s-1)\pi^{-s/2}\Gamma(s/2)\zeta(s)=2\xi(s)=g(s).
\end{equation}
The convention for $\xi$ is DLMF 25.4.4 \cite{zeta}.
Rapid decrease at both ends extends the left side to an entire
function; the identity continues to every $s$.

Set $G(z)=g(1/2+z)$. Substitution $x=e^t$ in (TS7), then evenness
of $K$, gives the absolutely convergent representation
\begin{equation}\tag{TS8}
 G(z)=2\int_0^\infty K(t)\cosh(zt)\dd t
 =\sum_{n=0}^\infty a_nz^{2n},\qquad
 a_n=\frac2{(2n)!}\int_0^\infty t^{2n}K(t)\dd t>0.
\end{equation}
Convergence is uniform on compact $z$-sets, by the exponential
majorants already proved. Thus these are the Taylor coefficients
of the original numerator.

\section{A fixed interval determined by actual zeros}
The residue-one pole of $\zeta$ at one, the Gamma recurrence, and
$\zeta(2)=\pi^2/6$, $\zeta(6)=\pi^6/945$ give
\begin{equation}\tag{TS9}
 G(1/2)=1,\qquad G(3/2)=\pi/3,\qquad G(11/2)=4\pi^3/63.
\end{equation}
These classical values and conventions are recorded in DLMF
25.2(i), 25.4.4 and 25.6.1 \cite{zeta,zetavalues}.
Since $(3/2)^{2n}\geq9(1/2)^{2n}$ for $n\geq1$, (TS8) gives
\[
 G(3/2)-a_0\geq9(G(1/2)-a_0).
\]
Subtracting the series at $3/2$ and $1/2$ also gives
$G(3/2)-G(1/2)\geq2a_1$. Thus
\begin{equation}\tag{TS10}
 a_0\geq\frac{27-\pi}{24}>\frac{167}{168},\qquad
 \frac{a_1}{a_0}\leq\frac{4(\pi-3)}{27-\pi}<\frac4{167}.
\end{equation}
The rational estimates use $\pi<22/7$, proved by a finite series
below. Positivity of all coefficients gives, for $|z|\leq11/2$,
\begin{equation}\tag{TS11}
 |G(z)|\geq a_0-|G(z)-a_0|
 \geq2a_0-G(11/2)
 >\frac{167}{84}-\frac4{63}\left(\frac{22}{7}\right)^3
 =\frac{1475}{86436}>0.
\end{equation}
Apply (TS11) to the actual zero $z=\delta+i\gamma$. It implies
$\delta^2+\gamma^2>121/4$ and hence
\begin{equation}\tag{TS12}
 \gamma^2-\delta^2>121/4-2(1/2)^2=119/4.
\end{equation}
No numerical zero table or extra hypothesis on the height is used.

The elementary inequalities $|\cos u|\leq1$ and
$\cos u\geq1-u^2/2$ in (TS8) give
$|G(iy)|\leq a_0$ and $G(iy)\geq a_0-a_1y^2$.
Put $T_*=27/5$ and $c_*=4/167$. The exact identities
\[
 119/4-T_*^2=59/100>0,\qquad
 1-c_*T_*^2=1259/4175>0
\]
show that squaring is legitimate on $0\leq y\leq T_*$:
\begin{equation}\tag{TS13}
 a_0^2(1-c_*y^2)^2\leq G(iy)^2\leq a_0^2.
\end{equation}

\section{The complete denominator and the entire outside mass}
Multiplying the four original linear factors on the source line gives
\begin{equation}\tag{TS14}\begin{aligned}
 H(y)&=[(y-\gamma)^2+\delta^2][(y+\gamma)^2+\delta^2]\\
 &=(y^2-\gamma^2+\delta^2)^2+4\delta^2\gamma^2>0,\\
 h(1/2+iy)&=H(y)^{m_0},\qquad
 w_h(y)=\frac{G(iy)^2}{2\pi H(y)^{2m_0}}.
\end{aligned}\end{equation}
In particular the full power and all four roots remain present.
Since $H'(y)=4y[y^2-(\gamma^2-\delta^2)]<0$ on $(0,T_*]$,
$R_h(y)=H(y)^{-2m_0}$ increases there.

For a finite positive measure $\nu$ and real integrable functions
$f,R$, direct expansion gives
\begin{equation}\tag{TS15}
 \left(\int fR\dd\nu\right)\left(\int1\dd\nu\right)
 -\left(\int f\dd\nu\right)\left(\int R\dd\nu\right)
 =\frac12\iint(f(y)-f(z))(R(y)-R(z))\dd\nu(y)\dd\nu(z).
\end{equation}
On $[0,T_*]$, use $d\nu=G(iy)^2dy$, $f(y)=\cosh(\alpha_*y)$
and $R=R_h$. These functions are bounded and the measure is finite,
so Fubini applies. The right side is nonnegative because both
$f$ and $R$ increase. All denominators below are positive by (TS13).
Using (TS13) only after this exact comparison yields
\begin{equation}\tag{TS16}
 \frac{\int_0^{T_*}\cosh(\alpha_*y)w_h(y)\dd y}
      {\int_0^{T_*}w_h(y)\dd y}
 \geq\frac{\int_0^{T_*}\cosh(\alpha_*y)G(iy)^2\dd y}
              {\int_0^{T_*}G(iy)^2\dd y}
 \geq J_*:=\frac1{T_*}\int_0^{T_*}\cosh(\alpha_*y)(1-c_*y^2)^2\dd y.
\end{equation}
On this interval $0<(1-c_*y^2)^2\leq1$, hence
$J_*\leq\cosh(\alpha_*T_*)$. On the entire remaining half-line
$y\geq T_*$, $\cosh(\alpha_*y)\geq\cosh(\alpha_*T_*)\geq J_*$.
Therefore
$\int_0^\infty(\cosh(\alpha_*y)-J_*)w_h(y)dy\geq0$:
the inner integral is nonnegative by (TS16) and the outer one
is nonnegative pointwise. Both integrals are finite by the retained
envelope. The exact evenness in (TS14) then gives
\begin{equation}\tag{TS17}
 \frac{M_h(\alpha_*)}{M_h(0)}
 =\frac{\int_0^\infty\cosh(\alpha_*y)w_h(y)\dd y}
        {\int_0^\infty w_h(y)\dd y}\geq J_*.
\end{equation}
This argument places no monotonicity requirement on $R_h$ outside
$[0,T_*]$ and removes no outer mass.

\section{Finite rational pressure estimate}
Because $\alpha_*>157/100$, the positive Taylor series of $\cosh$
and exact polynomial integration give
\begin{equation}\tag{TS18}\begin{aligned}
 J_*&\geq S_*:=\sum_{j=0}^8\frac{(157/100)^{2j}}{(2j)!}
 \left[\frac{T_*^{2j}}{2j+1}-\frac{2c_*T_*^{2j+2}}{2j+3}
                +\frac{c_*^2T_*^{2j+4}}{2j+5}\right],\\
 \frac{611585}{10000}&<S_*<\frac{611586}{10000},\qquad S_*>61.
\end{aligned}\end{equation}
Every bracket is exactly
$T_*^{-1}\int_0^{T_*}y^{2j}(1-c_*y^2)^2dy>0$.
For $x=41/10$, all terms after $x^{21}/21!$ have successive
ratios at most $x/22<1$, so
\begin{equation}\tag{TS19}
 e^{41/10}\leq\sum_{j=0}^{20}\frac{(41/10)^j}{j!}
 +\frac{(41/10)^{21}}{21!}\frac1{1-(41/10)/22}<61.
\end{equation}
Both (TS18) and (TS19) are finite rational inequalities verified
by the delivered exact cross-products. They imply
$\log(M_h(\alpha_*)/M_h(0))>41/10$. The positive $C^1$ function
$M_h$ satisfies the fundamental theorem of calculus on the whole
closed interval. Cauchy--Schwarz thus proves
\begin{equation}\tag{TS20}
 I_h\geq\frac1{\alpha_*}
 \left(\int_0^{\alpha_*}\frac{M_h'(\theta)}{M_h(\theta)}\dd\theta\right)^2
 =\frac2\pi\log^2\frac{M_h(\pi/2)}{M_h(0)}
 >\frac{1681}{50\pi}.
\end{equation}

\section{Physical EIQ measure and an independent exact moment derivation}
\label{eqphysical}
Use $a,b>0$ for the potential parameters in this section, to keep them
distinct from $\alpha_*=\pi/2$ and all original packet parameters.
Define the physical complete elliptic integrals
\[
 K(\kappa)=\int_0^{\pi/2}\frac{d\theta}{\sqrt{1-\kappa^2\sin^2\theta}},
 \qquad E(\kappa)=\int_0^{\pi/2}\sqrt{1-\kappa^2\sin^2\theta}\dd\theta.
\]
The conventions agree with Carlson's DLMF Chapter 19 \cite{elliptic}.
For $r\in(0,1)$, put $\kappa=\sqrt{1-r^2}$. Differentiating these
integrals gives $dE/dr=r(K-E)/(1-r^2)$ and
$dK/dr=(r^2K-E)/(r(1-r^2))$; the latter follows also by integrating
the derivative of $\sin\theta\cos\theta/
\sqrt{1-\kappa^2\sin^2\theta}$ between its zero boundary values.
It follows that
\begin{equation}\tag{TS21}
 \frac d{dr}\frac{rK}{E}
 =\frac{(K-E)(E-r^2K)}{(1-r^2)E^2}>0.
\end{equation}
The first numerator factor is the integral of
$(1-r^2)\sin^2\theta/\sqrt{1-(1-r^2)\sin^2\theta}$;
the second is the corresponding integral with $\cos^2\theta$.
Both are positive. Dominated convergence gives $rK\to0$, $E\to1$
as $r\downarrow0$, since the integrand of $rK$ is bounded by one.
At $r\uparrow1$, $K,E\to\pi/2$. Thus there is a unique solution of
\begin{equation}\tag{TS22}
 \frac{rK}{E}=\frac a{a+2},\qquad
 v=\frac{(a+2)\pi}{bE},\quad u=rv,
 \quad uK=\frac{a\pi}{b},\quad vE=\frac{(a+2)\pi}{b}.
\end{equation}

Put $A=u^2$, $B=v^2$, $m=(A+B)/2$, $R_s=\sqrt{(A+s)(B+s)}$,
$R_0=uv$, and $V(x)=b\sqrt x-a\log x$.
The substitutions $x=A\cos^2\theta+B\sin^2\theta$ give
\[
 \int_A^B\frac{V'(x)}{\sqrt{(B-x)(x-A)}}\dd x=0,\qquad
 \int_A^B\frac{xV'(x)}{\sqrt{(B-x)(x-A)}}\dd x=2\pi.
\]
Indeed the integrals with numerators $x^{-1/2},x^{-1},x^{1/2}$
are respectively $2K/v$, $\pi/(uv)$, and $2vE$.
Since $\int_0^\infty s^{-1/2}(x+s)^{-1}ds=\pi/\sqrt x$,
Tonelli in the first identity gives
$a/R_0=(b/(2\pi))\int_0^\infty ds/(\sqrt sR_s)$.
Define the positive function and measure
\begin{equation}\tag{TS23}\begin{aligned}
 H_*(x)&=\frac{a}{xR_0}-\frac b{2\pi}\int_0^\infty
           \frac{ds}{\sqrt s(x+s)R_s}
 =\frac b{2\pi x}\int_0^\infty\frac{\sqrt s}{(x+s)R_s}\dd s>0,\\
 d\rho_{a,b}(x)&=\frac{\sqrt{(B-x)(x-A)}}{2\pi}H_*(x)
                 \mathbf1_{(A,B)}(x)\dd x.
\end{aligned}\end{equation}
All auxiliary integrals converge at zero and infinity. This is exactly
the EIQ density \cite{EIQ}, with no coordinate rescaling.

For completeness its mass and Euler equality can be checked directly.
Let $R(z)=\sqrt{(z-A)(z-B)}$ with $R(z)/z\to1$ at infinity.
Substitution $x=m+(B-A)\cos\theta/2$ and then
$t=\tan(\theta/2)$ gives the arcsine transform $1/R(z)$.
Polynomial division of $(B-x)(x-A)$ then gives
\begin{equation}\tag{TS24}
 \frac1\pi\int_A^B\frac{\sqrt{(B-x)(x-A)}}{(x+s)(z-x)}\dd x
 =1-\frac{R_s+R(z)}{z+s}.
\end{equation}
The removable value at $z=-s$ is understood. The same division gives
$\pi^{-1}\int_A^B\sqrt{(B-x)(x-A)}(x+s)^{-1}dx=m+s-R_s$.
Using (TS23) and $\int_0^\infty s^{-1/2}(1-s/R_s)ds=2vE$,
which follows by averaging the Stieltjes representation of $\sqrt x$
against the arcsine measure, yields
\[
 \int d\rho_{a,b}=-a/2+bvE/(2\pi)=1.
\]
Furthermore (TS24), the Stieltjes identity, and the exact cancellation
$a/R_0=(b/(2\pi))\int_0^\infty ds/(\sqrt sR_s)$ give
\begin{equation}\tag{TS25}
 \int\frac{d\rho_{a,b}(x)}{z-x}
 =\frac{V'(z)-R(z)H_*(z)}2.
\end{equation}
These identities first hold off both $[A,B]$ and the negative real
axis, where all integrations are absolutely convergent. Taking real
boundary values on $(A,B)$ gives the principal-value derivative
of the logarithmic potential, because the density is locally $C^1$;
subtracting its value at the singular point proves this limit.
As $R(x+i0)$ is purely imaginary there, (TS25) implies
\begin{equation}\tag{TS26}
 V(x)-2\int\log|x-t|\dd\rho_{a,b}(t)=\ell\qquad(A\leq x\leq B).
\end{equation}
The endpoints follow by continuity; bounded compactly supported
density makes the logarithmic singularity uniformly integrable.
Outside $[A,B]$ its derivative is $R(x)H_*(x)$, negative on $(0,A)$
and positive on $(B,\infty)$. This proves the strict exterior Euler
inequality. It identifies the constructed measure with the unique
EIQ equilibrium: for the difference $\nu$ from any competing finite
energy probability measure, expansion of energy gives
$\int(V-2U_\rho-\ell)\dd\nu-\iint\log|x-y|\dd\nu(x)\dd\nu(y)$.
The first term equals its integral against the competing probability
and is nonnegative. The second is nonnegative, since for a zero-mass
signed measure the identity
\[
 -\iint\log|x-y|\dd\nu(x)\dd\nu(y)
 =\frac12\int_0^\infty\frac1t
       \iint e^{-t(x-y)^2}\dd\nu(x)\dd\nu(y)\dd t
\]
follows from the scalar identity for $-\log d$ by bounded truncation
and dominated convergence. The inner double integral is
$(4\pi t)^{-1/2}\int e^{-\zeta^2/(4t)}|\widehat\nu(\zeta)|^2d\zeta$.
It is positive unless $\nu=0$: a zero Fourier transform forces every
Gaussian convolution, and hence the measure, to vanish. The
logarithmic integrability needed here follows from finite energy and
the coercivity of $V$; cross terms with $\rho$ are integrable by its
bounded compact density. Infinite-energy competitors cannot improve
the finite energy of (TS23).

We now derive the closed logarithmic moment without importing ECL14.
Write
\[
 C=\frac{(u+v)^2}{4},\quad\eta=\frac{v-u}{v+u},\quad
 d\nu_0(x)=\frac{dx}{\pi\sqrt{(B-x)(x-A)}},\quad
 d\omega_0(x)=\frac{R_0}{x}\dd\nu_0(x).
\]
Both measures are probabilities. Under $x=m+(B-A)\cos\theta/2$
the first is $d\theta/\pi$ and
$x=C(1+2\eta\cos\theta+\eta^2)$,
$R_0=C(1-\eta^2)$. Therefore the second measure has the Poisson
kernel density $(1-\eta^2)/(1+2\eta\cos\theta+\eta^2)$ and
\begin{equation}\tag{TS27}
 \int\cos(j\theta)\dd\omega_0=(-\eta)^j\quad(j\geq0).
\end{equation}
This follows by expanding the two geometric series and integrating;
all series converge absolutely for $\eta<1$.
The absolutely convergent expansion of $\log x$ gives
\begin{equation}\tag{TS28}
 \int\log x\dd\nu_0=\log C,\qquad
 \int\log x\dd\omega_0=\log C+2\log(1-\eta^2)
 =2\log R_0-\log C.
\end{equation}
For $t=m+(B-A)\cos\phi/2\in[A,B]$, the logarithmic distance has
the Abel-convergent cosine expansion
\[
 \log|x-t|=\log\frac{B-A}{4}
       -2\sum_{j\geq1}\frac{\cos(j\theta)\cos(j\phi)}j.
\]
To verify it, factor $\cos\theta-\cos\phi$ into the two sine
factors and take real parts of $-\sum z^j/j=\log(1-z)$.
Inserting an Abel factor $s^j$, $0<s<1$, permits termwise
integration. As $s\uparrow1$, convergence holds in $L^1(d\theta)$
by the locally integrable logarithmic singularity; the Poisson
kernel in (TS27) is bounded, so the same is true against $d\omega_0$.
The endpoint values follow by this identical one-sided integrability.
Thus (TS27)--(TS28) give the exact potential identities
\begin{equation}\tag{TS29}
 \int\log|x-t|\dd\nu_0(x)=\log\frac{B-A}{4},\qquad
 \int\log|x-t|\dd\omega_0(x)=\log t+\log\eta.
\end{equation}
Finally the endpoint integrals in (TS22) give
\[
 \int\sqrt x\dd\nu_0=2vE/\pi=2(a+2)/b,\qquad
 \int\sqrt x\dd\omega_0=2uK/\pi=2a/b.
\]
Integrate (TS26) first against $\nu_0$, then against $\omega_0$.
Fubini applies because both auxiliary densities have integrable
square-root endpoints and the density of $\rho$ is bounded;
the logarithmic singularities are absolutely integrable. This yields
\begin{align*}
 \ell&=2(a+2)-a\log C-2\log((B-A)/4),\\
 2\mathcal L(a,b)+2\log\eta
 &=2a-a(2\log R_0-\log C)-\ell.
\end{align*}
Since $(B-A)/(4\eta)=C$, substitution proves the exact formula
\begin{equation}\tag{TS30}
 \boxed{\mathcal L(a,b)=(a+1)\log C-a\log R_0-2.}
\end{equation}
In particular, at $(a,b)=(2,\pi)$,
\begin{equation}\tag{TS31}
 uK=2,\quad vE=4,\quad 2rK=E,\qquad
 \mathcal L(2,\pi)=6\log((u+v)/2)-2\log(uv)-2.
\end{equation}
All these variables refer to the original physical potential.
Inserting $u=rv$ and $v=4/E$ in (TS3) gives
\begin{equation}\tag{TS32}
 \boxed{\mathcal C_\Gamma=4\log H_*,\qquad
 H_*:=\frac{(1+r)^3}{8rE(\sqrt{1-r^2})}.}
\end{equation}
In particular the denominator contains $E$, not $2E/\pi$.

\section{Exact finite bounds for the physical elliptic constant}
For rational $x\in(0,1)$ set $m_x=1-x^2$ and
\[
 t_n(x)=\left(\frac{\binom{2n}{n}}{4^n}\right)^2m_x^n,
 \quad K_N(x)=\sum_{n=0}^Nt_n(x),\quad
 E_N(x)=1-\sum_{n=1}^N\frac{t_n(x)}{2n-1}.
\]
The binomial series for $(1-z)^{-1/2}$, integrated against
$z=m_x\sin^2\theta$, gives $2K/\pi=\sum_{n\geq0}t_n$.
Here $\int_0^{\pi/2}\sin^{2n}\theta\,d\theta=
(\pi/2)\binom{2n}{n}/4^n$, obtained by integration by parts from
$n=0$. The binomial series for $(1-z)^{1/2}$ similarly gives
$2E/\pi=1-\sum_{n\geq1}t_n/(2n-1)$.
Since $t_{n+1}/t_n=m_x(2n+1)^2/(2n+2)^2<m_x$, geometric
majoration proves both finite enclosures
\begin{equation}\tag{TS33}
 K_N\leq\frac{2K}{\pi}\leq K_N+\frac{t_{N+1}}{1-m_x},\qquad
 E_N-\frac{t_{N+1}}{(2N+1)(1-m_x)}
 \leq\frac{2E}{\pi}\leq E_N.
\end{equation}
Take $N=512$, $r_-=1601/10000$, $r_+=801/5000$.
With each $K,E$ evaluated at $\sqrt{1-r_\pm^2}$, exact rational
substitution in (TS33) proves
\begin{equation}\tag{TS34}
 -\frac{443}{10^8}<2r_-\frac{2K_-}{\pi}-\frac{2E_-}{\pi}
 <-\frac{441}{10^8}<0,
 \quad
 0<\frac{2597}{10^7}<2r_+\frac{2K_+}{\pi}-\frac{2E_+}{\pi}
 <\frac{2598}{10^7}.
\end{equation}
The strict monotonicity (TS21) implies $r_-<r<r_+$.
The defining integrand shows that $E(\sqrt{1-r^2})$ increases
with $r$. Define $\underline E_-=E_{512}(r_-)-
t_{513}(r_-)/(1025(1-m_{r_-}))$. Using the physical $\pi/2$
factor and $157/50<\pi<22/7$, the following rational bounds hold:
\begin{equation}\tag{TS35}
 \frac{1176}{1000}
 <\frac{(1+r_-)^3}{8r_+(11/7)E_{512}(r_+)}
 <H_*<\frac{(1+r_+)^3}{8r_-(157/100)\underline E_-}
 <\frac{1179}{1000}<\frac{85}{72}<e^{1/6}.
\end{equation}
The penultimate inequality is rational. The last follows by retaining
the first three positive terms $1+1/6+1/72$ in the exponential series.
The left endpoint exceeds one. Equations (TS32) and (TS35) prove
$0<\mathcal C_\Gamma<2/3$.

For completeness all elementary constants used above have rational
certificates as well. The tangent addition formulas yield
$\tan(2\arctan(1/5))=5/12$,
$\tan(4\arctan(1/5))=120/119$, and
$\tan(4\arctan(1/5)-\arctan(1/239))=1$.
The final angle is positive, by
$\arctan x\geq x/(1+x^2)$ and $\arctan x\leq x$;
it is less than $4/5<1<\pi/2$, where
$\pi/2=2\arctan1>1$ follows by integrating $(1+t^2)^{-1}>1/2$
on $[0,1)$. Hence Machin's identity is
$\pi=16\arctan(1/5)-4\arctan(1/239)$.
The alternating series through eight terms, with its first omitted
term, encloses $\pi$ strictly between $157/50$ and $22/7$.
Also integration of the geometric series on $[0,1/3]$ gives
$\log2=2\sum_{j\geq0}(1/3)^{2j+1}/(2j+1)>56/81>69/100$.
All stated finite comparisons are made with exact integer
cross-multiplication by the delivered certificate; its 33 comparisons
pass identically with and without Python optimization.

\section{Coefficient and finite receivers on the original domain}
Combining (TS20), $\mathcal C_\Gamma<2/3$, $\log2>69/100$ and
$\pi<22/7$ gives
\begin{equation}\tag{TS36}
 -2\mathcal C_\Gamma+\frac{\log2}{\pi}I_h
 >-\frac43+\frac{1681}{50}\frac{69}{100}\frac{49}{484}
 =\frac{7370383}{7260000}>1.
\end{equation}
This completes the source and equilibrium part of Theorem~\ref{th:sign}.

The scalar receiving map is the already accepted original one
\cite{SBV}. To make its scope explicit, retain
\begin{equation}\tag{TS37}\begin{aligned}
 k&=4\ell+1\geq5,\quad e_k=1+k(m_0-1),\quad q=e_k(k+1)^2,\\
 \chi_k(S)&=\prod_{a,b=0}^k
 [S-k/2-(2a-k)\delta-i(2b-k)\gamma]^{e_k},\\
 G_N[\nu]&=(J_NH_N[\nu]^{-1}J_N^*)^{-1},\qquad J_N=\operatorname{rem}_{\chi_k},\\
 \mathcal B_k[\nu]&=\log\det G_{q-1}[\nu]+\log\det G_q[\nu]
 -\log\det G_{2q-1}[\nu]-\log\det G_{2q}[\nu],\\
 d\sigma(y)&=\frac{|\Gamma(1/4+iy/2)|^2}{2\pi}\dd y,\qquad
 \delta_{k,1}^{\rm ar}=\mathcal B_k[w_h^{*k}(y)dy]-\mathcal B_k[d\sigma].
\end{aligned}\end{equation}
Here $H_N$ is the original degree-$N$ polynomial Gram matrix in
the physical coordinate $S=k/2+iy$. The coordinate morphism is
$P(S)\mapsto P(k/2+iy)$, with inverse substitution
$y=(S-k/2)/i$; no roots or lower components are removed.
The Gamma mass is $\sqrt{2\pi}$ and the arithmetic mass is
$M_h(0)^k$. The accepted theorem states, on its fixed-packet
sufficiently-large domain $k\geq k_0(h)$,
\begin{equation}\tag{TS38}\begin{aligned}
 \delta_{k,1}^{\rm ar}
 &=-(4m_0-2)\mathcal C_\Gamma q+\frac{\log2}{\pi}I_h k^2+\epsilon_{h,k},\\
 |\epsilon_{h,k}|&\leq E_{h,k}=C_h^{\rm SBV}q\mathfrak r(q),\qquad
 \mathfrak r(q)=\frac{\log\log(q+16)}{\log(q+2)}
                       +q^{-1/200}\log^2(q+2)\longrightarrow0.
\end{aligned}\end{equation}
Its full accepted proof and providers are preserved separately.
The present argument establishes the missing sign by applying
that theorem, and does not reassign its proof as unfinished work.

For $m_0=1$, define $A_*=5683461/2420000$ and note
$A_*-4/3=\kappa_*$. Then $q=(k+1)^2$ and the exact finite lower
receiver is
\begin{equation}\tag{TS39}
 \delta_{k,1}^{\rm ar}
 >A_*k^2-\frac43q-E_{h,k}
 =\kappa_*q-A_*(2k+1)-E_{h,k}.
\end{equation}
Thus $\delta_{k,1}^{\rm ar}>q$ whenever $k\geq k_0(h)$ and
\begin{equation}\tag{TS40}
 C_h^{\rm SBV}\mathfrak r(q)+A_*\frac{2k+1}{q}<\kappa_*-1.
\end{equation}
This guard holds eventually for each fixed actual simple packet;
no packet-independent value for $k_0(h)$ or $C_h^{\rm SBV}$ is claimed.
Dividing (TS38) by $q$ gives the positive simple coefficient in (TS5).
At fixed $m_0\geq2$, $k^2/q\to0$, so it instead gives
$\delta_{k,1}^{\rm ar}/q\to-(4m_0-2)\mathcal C_\Gamma<0$.
The same pressure estimate remains true, but its $k^2$ term is
lower order on this original higher-multiplicity scale.

The exact four-endpoint receiving identity for any invertible marked
frame $\mathcal J$ is
$\det(\mathcal J^*G_N[\nu]\mathcal J)=|\det\mathcal J|^2\det G_N[\nu]$.
It follows by multiplicativity of determinants. If the same frame is
used at the four endpoints, its factor cancels with coefficients
$1+1-1-1=0$. The arithmetic-to-Gamma ratio of these complete
four-endpoint products consequently has logarithm
$\delta_{k,1}^{\rm ar}$ exactly; it exceeds $e^q$ under (TS40).
This proves the full-frame scalar transport without assigning a
sign to any proper projected subframe or individual pairing.

Finally the complete action receiver retains its existing identity
\begin{equation}\tag{TS41}
 J_k^{\rm action}
 =\mathcal B_k[d\sigma]+\delta_{k,1}^{\rm ar}
   +W_k^{\rm mono}-P_{k,-}-P_{k,+}.
\end{equation}
Substitution of (TS39) increases this receiver by the displayed
positive scalar correction on the simple stratum, while retaining
the full Gamma return, monic window, and both penalties. The
order-$k$ identity
$\delta_{k,k}^{\rm ar}=\delta_{k,1}^{\rm ar}
 -(\mathcal B_k^{(k)}-\mathcal B_k^{(1)})$ remains unchanged.
The seven absolute allocations and the projected responses $U,V$
are not evaluated by these scalar identities. Nothing here closes
the RH programme or infers a sign for the complete action after
omitting its other terms.

\begin{thebibliography}{9}
\bibitem{theta} W. P. Reinhardt and P. L. Walker,
\emph{Theta Functions}, NIST Digital Library of Mathematical Functions,
Chapter 20, equation 20.7.32,
\url{https://dlmf.nist.gov/20.7.E32}. Modular identity checked 19 September 2026.
\bibitem{zeta} T. M. Apostol, \emph{Zeta and Related Functions}, NIST DLMF,
Chapter 25, 25.2(i), 25.4.3--4,
\url{https://dlmf.nist.gov/25.2}, \url{https://dlmf.nist.gov/25.4}.
Residue and exact $\xi$ convention checked 19 September 2026.
\bibitem{zetavalues} T. M. Apostol, NIST DLMF, equation 25.6.1,
\url{https://dlmf.nist.gov/25.6.E1}. The displayed values cite
W. Magnus, F. Oberhettinger and R. P. Soni,
\emph{Formulas and Theorems for the Special Functions of Mathematical Physics},
1966, p. 19. Values checked 19 September 2026.
\bibitem{elliptic} B. C. Carlson, \emph{Elliptic Integrals}, NIST DLMF,
Ch. 19, definitions and derivatives, \url{https://dlmf.nist.gov/19.2},
\url{https://dlmf.nist.gov/19.4}. Physical integrals are defined explicitly
above, and the derivatives and rational tail bounds are derived here.
\bibitem{EIQ} Split-Zero programme, \emph{Exact equilibrium for the
square-root logarithmic potential}, EIQ1--28 and GEL providers,
\texttt{EIQ\_GEL\_original\_equilibrium.tex}, supplied and preserved source.
The physical density, endpoint equations, Euler equality and the
additional closed moment derivation are all written out above.
\bibitem{incoming} Supplied continuation, \emph{Positive original-source
pressure and the sign of the simple-packet arithmetic correction},
\texttt{ARITHMETIC\_SIGN\_FROM\_THETA.pdf}, 19 September 2026, nine pages,
APS1--36, and the corresponding supplied transcript. The pressure,
constant and sign assertions are independently rederived here.
\bibitem{SBV} Split-Zero programme, \emph{Signed arithmetic integrated
calculation}, \texttt{SIGNED\_ARITHMETIC\_INTEGRATED.tex}, accepted
19 September 2026, together with \texttt{CENTRAL\_FINITE\_REPAIR.tex},
\texttt{OUTER\_VALUE\_RECONSTRUCTION.tex}, the full localization, source-window,
original arithmetic envelope, and Gamma providers. The exact input
identities and hashes are listed in the accompanying acceptance record;
the sealed 243-page reader is preserved.
\end{thebibliography}
\end{document}
